
\begin{table}[!htb]
\centering
\caption{Percentage of seroconverted COVID-19 patients by time since symptom onset (Tan et al. \cite{Tan2020}). A total 312 tests were performed on 65 patients. Day is the number of days passed since COVID-19 symptom onset, Patients are the number of patients tested and IgG positive are the number of patients who tested positive for SARS-CoV-2 IgG antibodies.}

 \begin{tabular}{|c | c | c | c |} 
 \hline
 Day & Patients & IgG positive & \% \\ [0.5ex] 
 \hline\hline
 7 & 58 & 2 & 3.4 \\ 
 \hline
 10 & 62 & 12 & 19.4 \\
 \hline
 14 & 61 & 31 & 50.8 \\
 \hline
 21 & 54 & 32 & 59.3 \\
 \hline
 28 & 35 & 26 & 74.3 \\ 
 \hline
 35 & 22 & 17 & 77.3 \\
 \hline
 42 & 15 & 13 & 86.7 \\
 \hline
 49 & 5 & 4 & 80.0  \\ [1ex] 
 \hline
\end{tabular}

\label{tab:Tandata}
\end{table}


\begin{table}[!htb]
\caption{COVID-19 cases and serology survey results during spring 2020. The column COVID-19 cases (cumulative) shows the cumulative number and cumulative incidence of COVID-19 cases by the end of each week (Period). Data are shown both for the extended capital region of Finland (HUS) and for the target population of the current study (Study). The column Serological surveys (weekly) shows the weekly number of blood samples from the serological survey participants (Samples), the weekly number and proportion of positive samples using the screening test (Sreening pos. (\%)) and the weekly number and proportion of positive samples using the confirmation test (Confirmation pos. (\%)). }
\centering
\begin{tabular}{l|l|l|l|l|l}
\hline
\multicolumn{1}{c|}{} & \multicolumn{2}{c|}{COVID-19 cases\textsuperscript{a} (cumulative)} & \multicolumn{3}{c}{Serological surveys\textsuperscript{b} (weekly)} \\
\cline{2-3} \cline{4-6}
Period & HUS (\%) & Study (\%) & Samples & Screening pos. (\%) & Confirmation pos. (\%)\\
\hline
\cellcolor{gray!6}{10.02.2020 – 16.02.2020} & \cellcolor{gray!6}{0–10} & \cellcolor{gray!6}{0–10} & \cellcolor{gray!6}{-} & \cellcolor{gray!6}{-} & \cellcolor{gray!6}{-}\\
\hline
17.02.2020 – 23.02.2020 & 0–10 & 0–10 & - & - & -\\
\hline
\cellcolor{gray!6}{24.02.2020 – 01.03.2020} & \cellcolor{gray!6}{0–10} & \cellcolor{gray!6}{0–10} & \cellcolor{gray!6}{-} & \cellcolor{gray!6}{-} & \cellcolor{gray!6}{-}\\
\hline
02.03.2020 – 08.03.2020 & 24 (0) & 20 (0) & - & - & -\\
\hline
\cellcolor{gray!6}{09.03.2020 – 15.03.2020} & \cellcolor{gray!6}{220 (0.01)} & \cellcolor{gray!6}{190 (0.02)} & \cellcolor{gray!6}{-} & \cellcolor{gray!6}{-} & \cellcolor{gray!6}{-}\\
\hline
16.03.2020 – 22.03.2020 & 611 (0.04) & 505 (0.05) & - & - & -\\
\hline
\cellcolor{gray!6}{23.03.2020 – 29.03.2020} & \cellcolor{gray!6}{965 (0.06)} & \cellcolor{gray!6}{737 (0.07)} & \cellcolor{gray!6}{-} & \cellcolor{gray!6}{-} & \cellcolor{gray!6}{-}\\
\hline
30.03.2020 – 05.04.2020 & 1578 (0.09) & 1030 (0.1) & - & - & -\\
\hline
\cellcolor{gray!6}{06.04.2020 – 12.04.2020} & \cellcolor{gray!6}{2212 (0.13)} & \cellcolor{gray!6}{1332 (0.13)} & \cellcolor{gray!6}{23} & \cellcolor{gray!6}{1 (4.35)} & \cellcolor{gray!6}{0 (0)}\\
\hline
13.04.2020 – 19.04.2020 & 2825 (0.16) & 1621 (0.16) & 339 & 8 (2.36) & 1 (0.29)\\
\hline
\cellcolor{gray!6}{20.04.2020 – 26.04.2020} & \cellcolor{gray!6}{3436 (0.2)} & \cellcolor{gray!6}{1895 (0.19)} & \cellcolor{gray!6}{465} & \cellcolor{gray!6}{13 (2.8)} & \cellcolor{gray!6}{2 (0.43)}\\
\hline
27.04.2020 – 03.05.2020 & 3965 (0.23) & 2138 (0.21) & 170 & 4 (2.35) & 2 (1.18)\\
\hline
\cellcolor{gray!6}{04.05.2020 – 10.05.2020} & \cellcolor{gray!6}{4466 (0.26)} & \cellcolor{gray!6}{2415 (0.24)} & \cellcolor{gray!6}{139} & \cellcolor{gray!6}{2 (1.44)} & \cellcolor{gray!6}{0 (0)}\\
\hline
11.05.2020 – 17.05.2020 & 4804 (0.28) & 2636 (0.26) & 88 & 2 (2.27) & 1 (1.14)\\
\hline
\cellcolor{gray!6}{18.05.2020 – 24.05.2020} & \cellcolor{gray!6}{4987 (0.29)} & \cellcolor{gray!6}{2747 (0.28)} & \cellcolor{gray!6}{47} & \cellcolor{gray!6}{0 (0)} & \cellcolor{gray!6}{0 (0)}\\
\hline
25.05.2020 – 31.05.2020 & 5118 (0.3) & 2825 (0.28) & 48 & 0 (0) & 0 (0)\\
\hline
\cellcolor{gray!6}{01.06.2020 – 07.06.2020} & \cellcolor{gray!6}{5200 (0.3)} & \cellcolor{gray!6}{2863 (0.29)} & \cellcolor{gray!6}{48} & \cellcolor{gray!6}{2 (4.17)} & \cellcolor{gray!6}{1 (2.08)}\\
\hline
08.06.2020 – 14.06.2020 & 5240 (0.31) & 2885 (0.29) & 44 & 1 (2.27) & 0 (0)\\
\hline
\cellcolor{gray!6}{15.06.2020 – 21.06.2020} & \cellcolor{gray!6}{5279 (0.31)} & \cellcolor{gray!6}{2899 (0.29)} & \cellcolor{gray!6}{23} & \cellcolor{gray!6}{0 (0)} & \cellcolor{gray!6}{0 (0)}\\
\hline
22.06.2020 – 28.06.2020 & 5315 (0.31) & 2915 (0.29) & 9 & 0 (0) & 0 (0)\\
\hline
\cellcolor{gray!6}{29.06.2020 – 04.07.2020} & \cellcolor{gray!6}{5347 (0.31)} & \cellcolor{gray!6}{2932 (0.29)} & \cellcolor{gray!6}{22} & \cellcolor{gray!6}{2 (9.09)} & \cellcolor{gray!6}{0 (0)}\\
\hline
All weeks & 5347 (0.31) & 2932 (0.29) & 1465 & 35 (2.39) & 7 (0.48)\\
\hline
\multicolumn{6}{l}{\textsuperscript{a} HUS: COVID-19 cases in the extended capital region of Finland; Study: COVID-19 cases in the target }\\
\multicolumn{6}{l}{population of the current study. Populations 1.72M and 1.00M, respectively.}\\
\multicolumn{6}{l}{\textsuperscript{b} Results from the serological surveys for the target population of the current study. Samples gives the}\\
\multicolumn{6}{l}{number of blood samples taken during each week. Screening pos. (\%) gives the weekly number and proportion}\\
\multicolumn{6}{l}{of samples where SARS-CoV-2 IgG antibodies were detected. Confirmation pos. (\%) gives the weekly number}\\
\multicolumn{6}{l}{and proportion of positive samples confirmed via a microneutralisation test.}\\
\end{tabular}

\label{tab:covidcase}
\end{table}


\begin{table}[!htb]
\caption{Estimated and projected seroprevalence ($\pi^{(0)}(t)$ and $\pi^{(1)}(t)$) and the underreporting ratio ($\Delta(t)$, see equation \ref{eq:ratio}) of SARS-CoV-2 infections during the study period. The seroprevalences are shown in percentage scale. Displayed are the posterior means along with 95\% credible intervals, derived from the 2.5\% and 97.5\% quantiles of the distributions.}
\centering
\begin{tabular}{l|l|l|l|l|l}
\hline
\multicolumn{1}{c|}{ } & \multicolumn{1}{c|}{COVID-19 cases} & \multicolumn{2}{c|}{Confirmation test} & \multicolumn{2}{c}{Screening test} \\
\cline{2-2} \cline{3-4} \cline{5-6}
date & $\pi^{(1)}(t)$ & $\pi^{(0)}(t)$ & $\Delta(t)$ & $\pi^{(0)}(t)$ & $\Delta(t)$ \\
\hline
\cellcolor{gray!6}{09.04.2020} & \cellcolor{gray!6}{0.055 (0.050–0.061)} & \cellcolor{gray!6}{0.49 (0.20–0.91)} & \cellcolor{gray!6}{8.92 (3.64–16.53)} & \cellcolor{gray!6}{0.30 (0.022–0.91)} & \cellcolor{gray!6}{5.49 (0.40–16.53)}\\
\hline
16.04.2020 & 0.080 (0.072–0.087) & 0.49 (0.20–0.89) & 6.14 (2.54–11.26) & 0.30 (0.022–0.90) & 3.79 (0.28–11.34)\\
\hline
\cellcolor{gray!6}{23.04.2020} & \cellcolor{gray!6}{0.106 (0.096–0.114)} & \cellcolor{gray!6}{0.49 (0.21–0.89)} & \cellcolor{gray!6}{4.68 (1.95– 8.53)} & \cellcolor{gray!6}{0.30 (0.023–0.91)} & \cellcolor{gray!6}{2.89 (0.21– 8.58)}\\
\hline
30.04.2020 & 0.132 (0.121–0.141) & 0.50 (0.21–0.91) & 3.81 (1.59– 6.95) & 0.31 (0.023–0.92) & 2.37 (0.17– 7.02)\\
\hline
\cellcolor{gray!6}{07.05.2020} & \cellcolor{gray!6}{0.158 (0.145–0.168)} & \cellcolor{gray!6}{0.51 (0.22–0.94)} & \cellcolor{gray!6}{3.27 (1.37– 6.00)} & \cellcolor{gray!6}{0.32 (0.024–0.96)} & \cellcolor{gray!6}{2.05 (0.15– 6.12)}\\
\hline
14.05.2020 & 0.184 (0.170–0.194) & 0.53 (0.22–0.99) & 2.90 (1.20– 5.43) & 0.34 (0.024–1.01) & 1.84 (0.13– 5.52)\\
\hline
\cellcolor{gray!6}{21.05.2020} & \cellcolor{gray!6}{0.209 (0.194–0.220)} & \cellcolor{gray!6}{0.56 (0.23–1.06)} & \cellcolor{gray!6}{2.67 (1.08– 5.11)} & \cellcolor{gray!6}{0.36 (0.025–1.09)} & \cellcolor{gray!6}{1.72 (0.12– 5.24)}\\
\hline
28.05.2020 & 0.230 (0.214–0.241) & 0.58 (0.23–1.16) & 2.54 (1.01– 5.04) & 0.39 (0.027–1.21) & 1.69 (0.12– 5.28)\\
\hline
\cellcolor{gray!6}{04.06.2020} & \cellcolor{gray!6}{0.246 (0.231–0.257)} & \cellcolor{gray!6}{0.62 (0.24–1.29)} & \cellcolor{gray!6}{2.51 (0.96– 5.27)} & \cellcolor{gray!6}{0.43 (0.028–1.41)} & \cellcolor{gray!6}{1.73 (0.11– 5.74)}\\
\hline
11.06.2020 & 0.259 (0.244–0.268) & 0.66 (0.24–1.49) & 2.56 (0.93– 5.75) & 0.48 (0.029–1.73) & 1.86 (0.11– 6.69)\\
\hline
\cellcolor{gray!6}{18.06.2020} & \cellcolor{gray!6}{0.268 (0.254–0.276)} & \cellcolor{gray!6}{0.71 (0.25–1.76)} & \cellcolor{gray!6}{2.67 (0.92– 6.57)} & \cellcolor{gray!6}{0.56 (0.031–2.26)} & \cellcolor{gray!6}{2.08 (0.11– 8.50)}\\
\hline
25.06.2020 & 0.274 (0.262–0.281) & 0.78 (0.25–2.16) & 2.86 (0.91– 7.91) & 0.67 (0.032–3.14) & 2.43 (0.12–11.53)\\
\hline
\cellcolor{gray!6}{02.07.2020} & \cellcolor{gray!6}{0.279 (0.268–0.285)} & \cellcolor{gray!6}{0.88 (0.25–2.73)} & \cellcolor{gray!6}{3.14 (0.91– 9.80)} & \cellcolor{gray!6}{0.83 (0.033–4.66)} & \cellcolor{gray!6}{2.97 (0.12–16.71)}\\
\hline
\end{tabular}
\label{tab:results}
\end{table}

